\section*{Заключение}
\addcontentsline{toc}{section}{Заключение}

В ходе выполнения практической работы был развёрнут и настроен DNS-сервер BIND9
на виртуальной машине \texttt{gateway} под управлением Ubuntu~20.04~LTS.
Созданы прямая зона \texttt{STUDENT.GROUP.local} и обратная зона
\texttt{\cfgLabStudentN.168.192.in-addr.arpa}, обеспечивающие разрешение имён
хостов локальной сети в IP-адреса и обратно.

Настройка DDNS с ключом \texttt{rndc-key} позволила автоматизировать регистрацию
клиентов в DNS при получении IP-адреса через DHCP-сервер.
Клиентские машины автоматически получают имена
\texttt{HOSTNAME.STUDENT.GROUP.local}.

Полученные навыки~--- администрирование BIND9, работа с файлами зон,
настройка форвардеров и DDNS~--- являются базовыми компетенциями сетевого
администратора и широко применяются при построении корпоративных сетей.
